\begin{table}[H]
\centering
\caption{\label{tab:pi0_estimate_FF}Aggregate Skill Estimate: Proportion of True Zero-Alpha Active Funds -- FF Subperiod}
\centering
\resizebox{\linewidth}{!}{
\begin{threeparttable}
\begin{tabular}[t]{>{\raggedright\arraybackslash}p{16em}rrr>{\raggedright\arraybackslash}p{14em}}
\toprule
Metric & Estimate & $N$ & $\lambda$ & Interpretation\\
\midrule
$\hat{\pi}_0$: Proportion of True Zero-Alpha Active Funds & 76.0\% & 805 & 0.5 & Heterogeneous skill; majority zero-alpha\\
\bottomrule
\end{tabular}
\begin{tablenotes}
\item The proportion of true zero-alpha funds ($\hat{\pi}_0$) is estimated following Storey (2002) and Barras, Scaillet and Wermers (2010, \textit{Journal of Finance}), using p-values from Newey-West $t$-tests on full-period Carhart (1997) four-factor alphas. The estimator is $\hat{\pi}_0 = |\{p_i > \lambda\}| \,/\, [N(1-\lambda)]$, where $\lambda = 0.5$ is the standard tuning parameter following Storey (2002). Funds with p-values exceeding $\lambda$ are unlikely to have non-zero true alpha; the density of p-values in $(\lambda, 1]$ provides a conservative estimate of the zero-alpha proportion, bounded above at 1. Sample: actively managed funds ($N = 805$), Jan 1995--Sep 2006 (Fama-French 2010 subperiod). Passive and Unknown-classified funds are excluded. Compare with Fama and French (2010, \textit{Journal of Finance}), Table III, which reports an analogous bootstrap-based assessment over Jan 1984--Sep 2006.
\end{tablenotes}
\end{threeparttable}
}
\end{table}
